\documentclass[10pt,a4paper]{amsart}
\usepackage[margin=19mm]{geometry}
\usepackage{amsmath,amssymb,mathtools}
\usepackage[scr=rsfs,cal=euler]{mathalfa}
\usepackage{xfrac}
\usepackage[unicode,hidelinks,bookmarks=false]{hyperref}
\newcommand{\ie}{i.e.\ }
\newcommand{\cf}{cf.\ }
\newcommand{\resp}{respectively\ }
\newcommand{\Subsection}{\subsection}
\providecommand{\nobreakdash}{\nobreak\hbox{-}}
\setlength{\emergencystretch}{2em}
\multlinegap0pt
\usepackage{amssymb}
\usepackage{xcolor}
\usepackage{enumerate}
\usepackage{bm}
\newtheorem{proposition}{Proposition}
\newcommand{\ceil}[1]{\lceil #1 \rceil}
\providecommand{\e}{\varepsilon}
\title{A simple shallow neural network for emulating the solution to singularly perturbed problems}
\author{Christos Xenophontos, Aayushman Raina}
\date{Source: arXiv:2609.05180v1 (CC BY 4.0)}
\begin{document}
\maketitle

\noindent\emph{Source and licence.} A simple shallow neural network for emulating the solution to singularly perturbed problems. Version arXiv:2609.05180v1 (CC BY 4.0), distributed by arXiv under the Creative Commons Attribution 4.0 International licence, http://creativecommons.org/licenses/by/4.0/, as recorded in the arXiv licence metadata for this exact version and shown on the abstract page https://arxiv.org/abs/2609.05180v1. Full licence notice: \texttt{licenses/arxiv-2609.05180-CC-BY-4.0.txt}.\par\smallskip

\noindent\textit{Editorial scope.} Counted unit: the article's proposition prop:NNbl (source0.tex lines 322--329). The article's complete original proof of it is delivered with it (source0.tex lines 331--357, one \texttt{proof} environment of 27 lines). No mathematics is written, repaired or re-derived: the statement and the proof are the article's own lines, in the article's own notation, and no retained line is substituted or re-flowed.  Retained uncounted as the source-local context the proof needs: lines 143--168; lines 208--210.  Nothing was omitted from any retained span, so the package carries no omission record.  The retained lines cite 2 works, whose entries are transcribed verbatim from the article's own bibliography source in the e-print and delivered with short editorial labels recorded in the item's reference mapping.  No retained line contains a figure environment, an image-inclusion command or drawing code, so no image is delivered and the package declares no asset dependency.  Editorial formatting only, in addition: an AMS \texttt{amsart} preamble that restates the article's own declarations and macros used by the retained lines, namely the environments \texttt{proposition} on separate counters, together with the standard packages those lines need.  The retained environments print in this document as Proposition 1 in order of appearance, whereas the article prints the same items with the numbering of its own full document.  The complete native source of the article as distributed in the arXiv e-print for this exact version is bundled unchanged as \texttt{sources/209418/source0.tex}.
%%%%%%%%%%%%%%%%%%%%%%%%%%%%%%%%%%%%%%%%%%%%%%%%%%%%%%%%%%%
\section{Second order SPPs and the regularity of their solution}
\label{sec:model}
%%%%%%%%%%%%%%%%%%%%%%%%%%%%%%%%%%%%%%%%%%%%%%%%%%%%%%%%%%%
We consider the boundary value problem (BVP): find $u \in C^2(I) \cap C(\bar{I})$ such that
\begin{eqnarray}
-\e_1 u''(x) + \e_2  b(x)u'(x) + c(x) u(x) &=& f(x) \; , \; x \in I = (0,1) \; , \label{eq:de} \\
u(0) =  u(1)&=& 0, \label{eq:bc}
\end{eqnarray}
%
where $\e_1, \e_2 \in (0,1]$ and $b(x), c(x), f(x)$ are given analytic functions on
$\overline{I}=[0,1]$, satisfying $\forall \;x \in \overline{I}$,
$$
b(x) \geq \underline{b}>0, \quad c(x) \geq \underline{c}>0, \quad c(x)-\frac{\varepsilon_2}{2} b^{\prime}(x) \geq \underline{\delta} >0,
$$
%
for some positive constants $\underline{b}, \underline{c}, \underline{\delta} \in \mathbb{R}$, as well as
%
\begin{equation}
\label{eq:analytic}
\Vert f^{(n)} \Vert_{L^{\infty}(I)} \leq C_f K^n_f n! \: , \: \Vert b^{(n)} \Vert_{L^{\infty}(I)} \leq C_b K^n_b n! 
 \: , \: \Vert c^{(n)} \Vert_{L^{\infty}(I)} \leq C_c K^n_c n!
\end{equation}
%
$\forall \; n \in \mathbb{N}_0$,  where $C_f, K_f, C_b, K_b, C_c, K_c$, are positive constants independent of $\e_1$ and  $\e_2$.
(The superscript in parentheses denotes differentiation of order $n$.)

\begin{equation}\label{decomp}
u(x) = u_S(x) + u_{BL}^{-}(x) + u_{BL}^{+}(x) + u_R(x),
\end{equation}

\begin{proposition}\label{prop:NNbl}
Let $p \in \mathbb{N}$ be given and assume $u$ is the solution to (\ref{eq:de})--(\ref{eq:bc}), with $b(x)=b>0, c(x) = c >0$. Assuming $f$ is analytic, there exists a $\tanh$ NN $\Phi^{p}_{{SPP}}$ with size $M(\Phi^{p}_{{SPP}})$ and length $L(\Phi^{p}_{{SPP}})$, such that for all $\e_1 \ll \e_2^2 \ll 1$, the following estimate holds
%
$$
\Vert u - R[\Phi^{p}_{SPP}] \Vert_{L^{\infty}(I)} \leq C \left( e^{-\beta p} + e^{-\tilde{\delta} / \varepsilon_2} \right),
$$
for some positive constants $C, \beta, \tilde{\delta}$. Moreover, $M(\Phi^p_{SPP}) = O(p)$ and $L(\Phi^{p}_{{SPP}})= \ceil{3 \beta p/2}$.
\end{proposition}

\begin{proof}
We utilize the decomposition (\ref{decomp}) and construct separate NNs for each component, except for the remainder which is already exponentially small.

For the smooth (analytic) part $u_S$, we use \cite[Cor. 5.8]{Mishra} to deduce the existence of a $\tanh$ NN, say $\Phi^{\tanh}_S$, such that 
$$
\left\Vert u_S - R\left[\Phi^{\tanh}_{S}\right] \right\Vert_{L^{\infty}(I)} \leq C e^{-\beta p},
$$
with $M(\Phi^{\tanh}_S) = O(p)$, and $L(\Phi^{\tanh}_S) = \ceil{3 \beta p/2}$.

Since the coefficients are constant, the leading layer profiles are explicitly known, hence the shallow NN $\Phi^{\tanh}_{\text{exp}}$ emulates them within a given tolerance. In fact,
%
$$
\left\Vert u^{\pm}_{BL} - R^{\pm}\left[\Phi^{\tanh}_{\text{exp}}\right] \right\Vert_{L^{\infty}(I)} \leq C e^{-\beta p},
$$
where $R^{\pm}$ denotes the realization of the network at each endpoint.

We combine the three subnetworks in parallel (see, e.g., Proposition 2.3 in \cite{OPS}), and we obtain $\Phi^p_{SPP}$. The shallower exponential subnetworks are extended by identity layers so that the parallelization has the depth of the deepest component. Then,
$$
\begin{aligned}
\left\|u-R\left[\Phi_{\mathrm{SPP}}^p\right]\right\|_{L^{\infty}(I)} \leq & \left\|u_S-R\left[\Phi^{\tanh}_S\right]\right\|_{L^{\infty}(I)}+\left\|u_{BL}^{-}-R^{-}\left[\Phi_{\text{exp}}^{\tanh}\right]\right\|_{L^{\infty}(I)} \\
& +\left\|u_{BL}^{+}-R^{+}\left[\Phi_{\text{exp}}^{\tanh}\right]\right\|_{L^{\infty}(I)}+\left\|u_R\right\|_{L^{\infty}(I)} \\
\leq & C\left(e^{-\beta p}+e^{-\tilde{\delta} / \varepsilon_2}\right) .
\end{aligned}
$$
The resulting NN has size $M(\Phi_{SPP}^p)$ equal to the size of the three subnetworks, i.e.~$M(\Phi_{SPP}^p)=O(p)$,
and depth $L(\Phi_{SPP}^p) = \ceil{3 \beta p/2}=O(p)$.
\end{proof}

\begin{thebibliography}{99}
\bibitem[Mishra]{Mishra} T.~de Ryck, S.~Lanthaler and S.~Mishra, \emph{On the approximation of functions by tanh neural networks}, Neural Networks, \textbf{143}, (2021) 732 -- 750.
\bibitem[OPS]{OPS} J.~A.~A.~Opschoor, P.~C.~Petersen and Ch.~Schwab, \emph{Deep ReLU networks and high-order finite element methods}, Analysis and Applications, \textbf{18}, (2020) 715--770.
\end{thebibliography}
\end{document}
